\documentclass[11pt,a4paper]{article}
\usepackage[utf8]{inputenc}
\usepackage[T1]{fontenc}
\usepackage[brazil]{babel}
\usepackage{geometry}
\geometry{margin=2.5cm}
\usepackage{xcolor}
\usepackage{enumitem}
\usepackage{hyperref}
\usepackage{pifont}
\usepackage{tcolorbox}

\definecolor{reviewercolor}{RGB}{0,51,102}
\definecolor{responsecolor}{RGB}{0,100,0}
\definecolor{changecolor}{RGB}{139,69,19}

\newtcolorbox{reviewerbox}{
  colback=blue!5!white,
  colframe=reviewercolor,
  fonttitle=\bfseries,
  title=Comentário do Revisor,
  boxrule=0.5pt,
  left=4pt, right=4pt, top=4pt, bottom=4pt
}

\newcommand{\response}[1]{\par\noindent\textcolor{responsecolor}{\textbf{Resposta:}} #1}
\newcommand{\change}[1]{\par\noindent\textcolor{changecolor}{\textbf{Alterações no manuscrito:}} \textit{#1}}

\title{Resposta Ponto a Ponto aos Comentários dos Revisores\\[0.3em]
\large Manuscrito: ``Arquitetura da Plataforma Avalia+Tec com Segurança Criptográfica para\\
Processos Seletivos Transparentes e Auditáveis''}
\author{L.D.V. Santos, C.V.S. Oliveira, D.P.V. Santos, T.C. Azevedo,\\
R.N. Araújo Filho, J.M. Santos, M.M. Fernandes}
\date{\today}

\begin{document}
\maketitle

\noindent Prezado Editor e Revisores,

\medskip
\noindent Agradecemos sinceramente ao Editor e a ambos os Revisores pela avaliação minuciosa, construtiva e perspicaz do nosso manuscrito. Seus comentários identificaram fragilidades importantes que melhoraram substancialmente o rigor científico, a transparência e a honestidade do artigo.

Abordamos todos os comentários levantados por ambos os revisores. A seguir, apresentamos nossas respostas ponto a ponto.

\bigskip
\hrule
\bigskip

%% ================================================================
\section*{Revisor \#1}
%% ================================================================

\subsection*{Comentário 1.1 --- Modelo de ameaças e alegação de ``imutabilidade''}

\begin{reviewerbox}
O modelo de ameaças contém uma vulnerabilidade crítica que os próprios autores reconhecem: um administrador de servidor com acesso root ao banco de dados poderia reescrever toda a cadeia de hashes de forma consistente. Isso compromete fundamentalmente a alegação central do artigo sobre uma trilha de auditoria ``imutável''. Sem ancoragem externa de hashes (por exemplo, blockchain público ou autoridade de carimbo de tempo compatível com RFC~3161), o uso do termo ``imutável'' é tecnicamente incorreto e enganoso. Os autores devem implementar essa salvaguarda ou revisar substancialmente suas alegações.
\end{reviewerbox}

\response{Concordamos plenamente com esta crítica. O termo ``imutável'' era tecnicamente impreciso na ausência de ancoragem externa, e agradecemos ao revisor por insistir nesse ponto.}

Na versão revisada, tomamos ambas as ações corretivas sugeridas pelo revisor:

\begin{enumerate}[nosep]
  \item Implementação de Carimbo de Tempo Confiável RFC~3161. Implementamos um módulo \texttt{TimestampService} que submete o resumo SHA-256 de cada inscrição a uma Autoridade de Carimbo de Tempo (TSA) externa compatível com RFC~3161 (padrão: FreeTSA.org). A TSA retorna um token TimeStampResp (TSR) assinado criptograficamente, que é armazenado junto ao registro da inscrição. Mesmo um administrador com acesso root não pode forjar um TSR válido sem a chave privada de assinatura da TSA. O código inclui codificação ASN.1 DER, requisição HTTP POST à TSA e lógica de verificação do TSR. Uma migração SQL correspondente (\texttt{002\_add\_rfc3161\_timestamp.sql}) adiciona as colunas \texttt{tsr\_base64}, \texttt{tsr\_tsa\_url} e \texttt{tsr\_requested\_at}. A rota \texttt{/api/verify/:protocol} agora retorna os resultados da verificação do TSR juntamente com a verificação de hash existente.

  \item Correção terminológica. Todas as ocorrências de ``imutável'' foram substituídas por ``à prova de adulteração'' em ambas as versões (inglês e português) do manuscrito, no resumo, nos destaques, no rótulo do diagrama de arquitetura e em todas as descrições técnicas.
\end{enumerate}

\change{%
Resumo reescrito (ambas as versões); Seção~1 motivação; Seção~2 Funcionalidades do Software --- novo parágrafo descrevendo RFC~3161; Seção~3 Modelo de Ameaças --- completamente reescrita com três classes de adversários (atacante externo, insider sem privilégios, administrador privilegiado) e análise de risco residual; diagrama de arquitetura: ``Armazenamento Imutável de Eventos'' $\to$ ``Armazenamento à Prova de Adulteração''; Seção~7 Conclusões --- linguagem de resistência a adulteração. Novos arquivos: \texttt{TimestampService.js}, \texttt{002\_add\_rfc3161\_timestamp.sql}. Arquivos modificados: \texttt{RegisterSubmission.js}, \texttt{SqlSubmissionRepository.js}, \texttt{JsonSubmissionRepository.js}, \texttt{index.js}.}

\bigskip

\subsection*{Comentário 1.2 --- Evidência empírica insuficiente}

\begin{reviewerbox}
A evidência empírica para os ganhos reportados de desempenho (redução de 65\% no tempo de ciclo, redução de 45,4\% no tempo do avaliador) é insuficiente. A amostra ($n=60$ inscrições, $n=6$ avaliadores) provém de uma única instituição, e as comparações com fluxos anteriores baseiam-se exclusivamente em dados retrospectivos de pesquisa administrativa auto-reportados, suscetíveis a viés de memória. Nenhum intervalo de confiança, teste formal de significância (ex: Mann--Whitney U) ou comparação com grupo controle é fornecido.
\end{reviewerbox}

\response{Concordamos que a apresentação estatística original era inadequada. O manuscrito revisado agora inclui:}

\begin{itemize}[nosep]
  \item Intervalos de confiança de 95\% para todas as métricas reportadas;
  \item Testes de Mann--Whitney $U$ (bicaudais) com $p$-valores exatos: tempo de ciclo $U=0$, $p=0{,}029$; tempo do avaliador $U=2$, $p=0{,}009$;
  \item Tamanhos de efeito: rank-biserial $r=1{,}0$ para tempo de ciclo, Cohen's $d=1{,}88$ (grande) para tempo do avaliador;
  \item Teste exato de Fisher para taxas de contestação: 0/60 vs.\ 14/292, $p=0{,}069$ (unilateral), explicitamente notado como \emph{não atingindo significância convencional};
  \item Reconhecimento explícito de viés de memória, efeito Hawthorne, tamanho amostral pequeno e limitação de instituição única;
  \item Linguagem moderada: ``estimativas preliminares de tamanho de efeito'', ``associado a'', ``sugere'' em vez de alegações causais.
\end{itemize}

Não adicionamos grupo controle, pois se tratou de uma implantação institucional real onde randomização seria impraticável e eticamente problemática. Essa limitação é explicitamente declarada (Seção~6).

\change{Seção~5 (Impacto) --- completamente reescrita com testes estatísticos e linguagem moderada. Seção~6 (Limitações) --- novas subseções ``Desenho do estudo'' e ``Qualidade dos dados de baseline''.}

\bigskip

\subsection*{Comentário 1.3 --- Novidade superestimada}

\begin{reviewerbox}
A novidade da contribuição criptográfica é superestimada. Hashing SHA-256 e autenticação HMAC-SHA256 em ambiente Node.js representam prática padrão de engenharia de software, não uma contribuição científica. A alegação de que nenhuma plataforma existente integra esses mecanismos requer justificativa mais forte do que a comparação apresentada na Tabela~3.
\end{reviewerbox}

\response{Concordamos plenamente. As primitivas criptográficas individuais (SHA-256, HMAC-SHA256) são de fato bem estabelecidas e sua implementação não constitui uma alegação de novidade científica. O manuscrito revisado reposiciona explicitamente a contribuição:}

\begin{itemize}[nosep]
  \item Um novo parágrafo ``Escopo da contribuição'' na Seção~1 declara: ``\textit{A contribuição científica reside nessa integração sistemática, e não nas primitivas criptográficas subjacentes.}''
  \item O resumo agora diz: ``\textit{A plataforma combina técnicas criptográficas bem estabelecidas---[\ldots]---em um fluxo de trabalho unificado e de código aberto especificamente projetado para [\ldots]}''
  \item Todas as alegações de novidade foram substituídas pelo enquadramento de ``engenharia de integração''.
\end{itemize}

\change{Resumo, Seção~1 (novo parágrafo ``Escopo da contribuição''), Seção~7 (Conclusões) --- contribuição explicitamente enquadrada como engenharia de integração.}

\bigskip

\subsection*{Comentário 1.4 --- Tabela comparativa superficial}

\begin{reviewerbox}
A análise comparativa na Tabela~3 é superficial e parece construída para favorecer o avalia+Tec. Notavelmente, o artigo alega conformidade com a LGPD como vantagem exclusiva sem avaliar os mecanismos de conformidade com o GDPR nas plataformas concorrentes. Os critérios de comparação devem ser definidos independentemente do conjunto de funcionalidades do sistema proposto.
\end{reviewerbox}

\response{Concordamos que o formato binário Sim/Não original era redutor e os critérios não eram suficientemente independentes. A Tabela~3 revisada incorpora:}

\begin{itemize}[nosep]
  \item Escala de pontuação em três níveis (\ding{108} completo/nativo, \ding{119} parcial/plugin, \ding{55} ausente) para evitar a simplificação binária;
  \item Duas novas colunas que não favorecem o avalia+Tec: ``TSA Ext.'' (avalia+Tec recebe \ding{119} parcial, pois a implementação é nova) e ``Escala Comprovada'' (avalia+Tec recebe \ding{55}, reconhecendo sua implantação em instituição única vs.\ sistemas em produção com milhões de submissões);
  \item Uma coluna ``GDPR/LGPD'' substituindo a coluna exclusiva de LGPD, com notas sobre conformidade via plugins (ex: controles de privacidade do OpenReview);
  \item Uma coluna ``Plugins Ext.'' creditando explicitamente HotCRP e OpenReview por seus ecossistemas de extensões;
  \item Notas metodológicas explicando os critérios de pontuação: ``avaliado via documentação publicada, inspeção de código-fonte e matrizes oficiais de funcionalidades.''
\end{itemize}

\change{Tabela~3 completamente reformulada (10 colunas, escala de 3 níveis, notas de rodapé, coluna GDPR, coluna ``Escala Comprovada'').}

\bigskip

\subsection*{Comentário 1.5 --- Categoria de submissão incorreta}

\begin{reviewerbox}
O manuscrito foi submetido na categoria ``Humanidades, Artes e Ciências Sociais'', o que é claramente incorreto para uma contribuição de engenharia de software. Os autores devem reenviar sob uma área temática apropriada.
\end{reviewerbox}

\response{Agradecemos ao revisor por apontar esse erro administrativo. A resubmissão foi reclassificada na área ``Ciência da Computação'', que é a categoria apropriada para uma contribuição de engenharia de software descrevendo uma plataforma web com componentes criptográficos. Pedimos desculpas pelo equívoco.}

\change{Metadados da submissão corrigidos para categoria ``Ciência da Computação''.}

\bigskip
\hrule
\bigskip

%% ================================================================
\section*{Revisor \#2}
%% ================================================================

\subsection*{Comentário 2.1 --- Engenharia de integração, não novidade científica}

\begin{reviewerbox}
O manuscrito, em grande parte, integra primitivas criptográficas bem estabelecidas (SHA-256, HMAC, RBAC, PostgreSQL RLS) sem propor inovações metodológicas, algorítmicas ou arquiteturais. A contribuição parece ser primariamente um esforço de implementação, e não um avanço cientificamente novo. Os autores devem articular claramente o que é fundamentalmente novo além da engenharia de integração.
\end{reviewerbox}

\response{Concordamos e reposicionamos explicitamente a contribuição como engenharia de integração. Veja nossa resposta ao Revisor~\#1, Comentário~1.3. O manuscrito revisado declara que a contribuição reside na \emph{integração sistemática} de primitivas estabelecidas em um fluxo de trabalho específico do domínio, e não nas primitivas em si. Também observamos que o escopo da SoftwareX abrange explicitamente ``software que tem impacto significativo na pesquisa'', independentemente de novidade algorítmica.}

\change{Mesmo que Comentário~1.3. Resumo, Seção~1 (parágrafo ``Escopo da contribuição''), Seção~7.}

\bigskip

\subsection*{Comentário 2.2 --- Estatística inferencial insuficiente}

\begin{reviewerbox}
A redução de 65\% no tempo de ciclo e o ganho de 45,4\% na eficiência do avaliador reportados baseiam-se em uma implantação pequena de instituição única ($n=60$ inscrições, $n=6$ avaliadores). Nenhuma estatística inferencial, intervalo de confiança ou teste de hipóteses é fornecido. Alegações sobre eliminação de disputas e melhorias de eficiência são, portanto, superestimadas e insuficientemente fundamentadas.
\end{reviewerbox}

\response{Integralmente abordado. Veja nossa resposta ao Revisor~\#1, Comentário~1.2. A Seção~5 revisada agora inclui ICs de 95\%, testes de Mann--Whitney $U$ com $p$-valores exatos, tamanhos de efeito e teste exato de Fisher para taxas de contestação. Todas as alegações foram adequadamente moderadas.}

\change{Mesmo que Comentário~1.2. Seção~5 (Impacto), Seção~6 (Limitações).}

\bigskip

\subsection*{Comentário 2.3 --- Dados retrospectivos auto-reportados}

\begin{reviewerbox}
A avaliação baseia-se em dados administrativos retrospectivos auto-reportados. Não há baseline controlado, comparação randomizada nem validação longitudinal multi-coorte. Isso enfraquece a credibilidade empírica das melhorias de desempenho reportadas.
\end{reviewerbox}

\response{Reconhecemos essa limitação explicitamente na Seção~6 revisada (Limitações), na nova subseção ``Qualidade dos dados de baseline''. O texto agora declara: ``\textit{As comparações de baseline baseiam-se em dados retrospectivos de pesquisa administrativa auto-reportados e estão sujeitas a viés de memória. Os avaliadores podem superestimar o tempo anteriormente gasto, inflando o ganho aparente de eficiência. Não existem dados objetivos de rastreamento de tempo para o período pré-implantação.}'' Também observamos a ausência de comparação randomizada e explicamos por que a randomização era impraticável em uma implantação institucional real. Validação longitudinal multi-institucional é identificada como trabalho futuro planejado.}

\change{Seção~5 --- alertas explícitos sobre viés de memória e efeito Hawthorne. Seção~6 --- subseção ``Qualidade dos dados de baseline''. Seção~7 --- trabalho futuro inclui implantações multi-institucionais com desenhos quase-experimentais.}

\bigskip

\subsection*{Comentário 2.4 --- Fragilidade do administrador confiável}

\begin{reviewerbox}
O manuscrito assume explicitamente um administrador confiável, o que enfraquece fundamentalmente as garantias de integridade. Uma vez que um usuário root malicioso pode reescrever toda a cadeia de hashes, o sistema não fornece resistência forte a adulterações em contextos institucionais adversariais. As alegações de segurança devem ser atenuadas em conformidade.
\end{reviewerbox}

\response{Abordado pela mesma implementação de RFC~3161 descrita na nossa resposta ao Revisor~\#1, Comentário~1.1. A Seção~3 revisada agora apresenta um modelo de ameaças com três classes:}

\begin{enumerate}[nosep]
  \item Atacante externo (não autenticado) --- mitigado por TLS, limitação de taxa, sanitização de entrada, autenticação JWT;
  \item Insider sem privilégios (usuário autenticado) --- mitigado por RBAC, RLS, formulários autenticados por HMAC, trilha de auditoria com cadeia de hashes;
  \item Administrador privilegiado (acesso root ao banco) --- mitigado por ancoragem externa de hashes via RFC~3161.
\end{enumerate}

\noindent A limitação residual é explicitamente declarada: se o administrador suprimir o armazenamento do TSR \emph{antes} da solicitação do carimbo de tempo, nenhuma âncora externa existirá para aquele registro. Isso é identificado como área para melhoria futura (ex: ancoragem contínua em blockchain, gerenciamento de chaves multipartidário).

\change{Seção~3 --- reescrita completa. Novo código: \texttt{TimestampService.js}. Todo ``imutável'' $\to$ ``à prova de adulteração''.}

\bigskip

\subsection*{Comentário 2.5 --- Comparação sem rigor}

\begin{reviewerbox}
A comparação com plataformas como SIGAA, OpenReview, EasyChair, HotCRP e SUCUPIRA carece de rigor metodológico. A comparação funcional não considera capacidades de plugins, extensibilidade ou possibilidades alternativas de integração criptográfica. Apresenta uma simplificação binária ``Sim/Não'' que pode deturpar os sistemas concorrentes.
\end{reviewerbox}

\response{Integralmente abordado. Veja nossa resposta ao Revisor~\#1, Comentário~1.4. A Tabela~3 revisada utiliza pontuação em três níveis, credita explicitamente ecossistemas de plugins, inclui colunas onde o avalia+Tec recebe nota inferior aos concorrentes e documenta a metodologia de comparação em notas de rodapé.}

\change{Mesmo que Comentário~1.4. Tabela~3 reformulada.}

\bigskip

\subsection*{Comentário 2.6 --- Sem validação formal de segurança}

\begin{reviewerbox}
Embora SHA-256 e HMAC sejam citados, não há validação formal de segurança, nenhum relatório de teste de penetração, nenhuma auditoria externa e nenhuma verificação de práticas corretas de gerenciamento de chaves. Simplesmente chamar bibliotecas-padrão não garante implantação segura.
\end{reviewerbox}

\response{Concordamos que chamar bibliotecas-padrão não garante implantação segura. O manuscrito revisado aborda isso de duas formas:}

\begin{enumerate}[nosep]
  \item Evidências positivas: A Seção~3 agora declara que a implementação foi verificada contra o checklist OWASP Top~10 e que o \texttt{npm audit} não reporta vulnerabilidades de alta severidade conhecidas nas dependências de produção. Também descrevemos as práticas atuais de gerenciamento de chaves (segredos JWT e chaves HMAC armazenados como variáveis de ambiente, nunca em código-fonte).
  \item Limitação honesta: As Seções~3 e~6 (subseção ``Segurança'') declaram explicitamente: ``\textit{Nenhum teste formal de penetração ou auditoria externa de segurança foi conduzido até o momento; isso representa trabalho futuro planejado.}'' e ``\textit{Esses não constituem uma avaliação abrangente de segurança.}''
\end{enumerate}

Teste formal de penetração e auditoria externa são listados como trabalho futuro na Seção~7.

\change{Seção~3 (Modelo de Ameaças) --- menção a OWASP/npm audit, descrição de gerenciamento de chaves. Seção~6 (Limitações) --- nova subseção ``Segurança''.}

\bigskip

\subsection*{Comentário 2.7 --- Sem teste de estresse}

\begin{reviewerbox}
Os testes de desempenho foram realizados com 60 inscrições. Nenhum teste de estresse, simulação de usuários concorrentes ou benchmarking de alta carga é apresentado. Não está claro como o sistema se comporta com centenas ou milhares de usuários simultâneos.
\end{reviewerbox}

\response{Agora conduzimos e reportamos testes de carga. Um novo script de teste de carga (\texttt{scripts/load-test.js}) utilizando autocannon (ferramenta de benchmarking HTTP para Node.js) foi criado com três cenários: GET~/ (página inicial), GET~/api/verify/:protocol (verificação de integridade) e POST~/api/inscricao (criação de inscrição com hashing SHA-256 e carimbo de tempo RFC~3161 assíncrono). Os testes foram executados contra uma instância do servidor avalia+Tec em operação, com 100 e 500 conexões simultâneas por 30~segundos cada.}

Uma nova Seção~4.1 (Testes de Carga) e Tabela~5 reportam os resultados medidos:
\begin{itemize}[nosep]
  \item A 100 conexões: GET~/ alcança 1.753~req/s (p50=52\,ms, p97,5=98\,ms); POST~/api/inscricao alcança 1.338~req/s (p50=68\,ms, p97,5=133\,ms);
  \item A 500 conexões: operações de leitura sustentam $>$1.700~req/s, mas a latência de cauda do POST aumenta acentuadamente (p99=8.153\,ms) com 1.127 erros de 47.976 requisições, indicando saturação do event loop que confirma a necessidade de escalonamento horizontal além de $\sim$500 escritores concorrentes;
  \item O texto declara explicitamente que escalonamento horizontal (balanceamento de carga, réplicas de leitura PostgreSQL) é uma extensão arquitetural direta para implantações maiores.
\end{itemize}

Os resultados brutos completos (JSON estruturado com todos os percentis de latência, throughput e contagem de erros para todas as seis combinações cenário/concorrência) estão incluídos no pacote de replicação (\texttt{load-test/load\_test\_results.json}).

\change{Nova Seção~4.1 (Testes de Carga), nova Tabela~5 com dados medidos. Novo arquivo: \texttt{scripts/load-test.js}. Pacote de replicação inclui resultados JSON brutos.}

\bigskip

\subsection*{Comentário 2.8 --- Conformidade com LGPD não demonstrada}

\begin{reviewerbox}
Alegações de conformidade com a LGPD são afirmadas, mas não rigorosamente demonstradas através de mapeamento formal de conformidade, DPIA (Avaliação de Impacto sobre a Proteção de Dados) ou referências de auditoria jurídica. As declarações de conformidade parecem descritivas e não validadas.
\end{reviewerbox}

\response{Reconhecemos que a plena conformidade com a LGPD requer medidas organizacionais, jurídicas e técnicas, e uma única plataforma de software não pode ``estar em conformidade com a LGPD'' isoladamente. O manuscrito revisado:}

\begin{enumerate}[nosep]
  \item Substitui ``em conformidade com a LGPD'' por ``design compatível com a LGPD'' em todo o texto;
  \item Adiciona uma nova Tabela~(tab:lgpdmapping) que mapeia artigos específicos da LGPD (Art.~6~I--II, Art.~6~V, Art.~6~VII--VIII, Art.~6~X, Art.~46, Art.~48, Art.~18) aos controles técnicos correspondentes implementados no avalia+Tec;
  \item Adiciona um disclaimer explícito: ``\textit{A plena conformidade com a LGPD requer políticas institucionais, um RIPD (Relatório de Impacto à Proteção de Dados Pessoais) formal e revisão jurídica por cada organização que implantar o sistema. Nenhum RIPD formal foi conduzido para a implantação atual.}''
\end{enumerate}

\change{Seção~2 --- terminologia ``design compatível com a LGPD'', nova Tabela~(tab:lgpdmapping), disclaimer. Seção~6 (Limitações) --- nova subseção ``Conformidade com a LGPD''. Seção~7 --- RIPD listado como trabalho futuro.}

\bigskip

\subsection*{Comentário 2.9 --- Sem dataset público ou pacote de replicação}

\begin{reviewerbox}
Embora métricas anonimizadas sejam ditas ``disponíveis mediante solicitação'', nenhum dataset público ou pacote de replicação é fornecido. Isso limita a transparência e a reprodutibilidade, que são centrais aos padrões da SoftwareX.
\end{reviewerbox}

\response{Concordamos. A declaração revisada de Disponibilidade de Dados agora diz:}

\begin{quote}
``\textit{Métricas anonimizadas de implantação, scripts de análise estatística (Python) e o framework de teste de carga estão disponíveis no pacote de replicação em \url{https://doi.org/10.5281/zenodo.18868994}.}''
\end{quote}

O pacote de replicação está publicado no Zenodo em \url{https://doi.org/10.5281/zenodo.18868994}. O pacote contém os seguintes arquivos de evidência:
\begin{itemize}[nosep]
  \item \texttt{data/deployment\_metrics.csv} --- métricas anonimizadas de implantação (sem PII);
  \item \texttt{analysis/statistical\_analysis.py} --- script Python reproduzindo todos os ICs, testes de Mann--Whitney $U$, tamanhos de efeito e teste exato de Fisher;
  \item \texttt{analysis/output\_statistical\_analysis.txt} --- saída capturada confirmando os valores do artigo;
  \item \texttt{load-test/load-test.js} --- script de benchmarking autocannon;
  \item \texttt{load-test/load\_test\_results.json} --- JSON estruturado com resultados brutos de todas as seis combinações cenário/concorrência.
\end{itemize}

Metadados de citação legíveis por máquina são fornecidos via \texttt{CITATION.cff} e \texttt{.zenodo.json} na raiz do repositório.

\change{Declaração de Disponibilidade de Dados atualizada. Seção~6 (Limitações) --- subseção ``Reprodutibilidade'' referencia Zenodo. Seção~7 --- pacote de replicação mencionado. \texttt{CITATION.cff} e \texttt{.zenodo.json} adicionados.}

\bigskip

\subsection*{Comentário 2.10 --- Manuscrito prolixo}

\begin{reviewerbox}
O artigo parece prolixo para uma publicação de software original na SoftwareX. Seções descrevendo fundamentos criptográficos e diagramas de fluxo poderiam ser condensadas. Algumas explicações conceituais são repetitivas e inflam o manuscrito sem adicionar profundidade técnica.
\end{reviewerbox}

\response{Condensamos o manuscrito da seguinte forma:}

\begin{itemize}[nosep]
  \item Removendo explicações redundantes de fundamentos de SHA-256 e HMAC (leitores da SoftwareX são considerados familiarizados com hashing padrão);
  \item Consolidando texto motivacional repetido entre seções;
  \item Compactando a discussão da comparação (a Tabela~3 redesenhada transmite mais informação em menos texto);
  \item Movendo descrições detalhadas de fluxo para o diagrama de arquitetura e sua legenda;
  \item Usando linguagem mais concisa e direta em todo o texto.
\end{itemize}

\noindent Apesar da adição de novo conteúdo técnico (implementação RFC~3161, testes de carga, tabela de mapeamento LGPD, modelo de ameaças expandido, testes estatísticos, limitações estruturadas), o texto geral está mais focado e menos repetitivo. Estimamos uma redução líquida no volume de prosa de aproximadamente 15\% (excluindo tabelas e equações).

\change{Ao longo de todo o manuscrito. Particularmente: Seção~1 (motivação condensada), Seção~2 (primer criptográfico removido), Seção~5 (tabelas substituem descrições prolixas).}

\bigskip
\hrule
\bigskip

\section*{Resumo das Alterações}

\begin{table}[h]
\centering
\small
\begin{tabular}{|c|p{5cm}|p{7cm}|}
\hline
\textbf{Rev./Pt} & \textbf{Questão} & \textbf{Ação Tomada} \\
\hline
R1.1 / R2.4 & Alegação de ``imutável''; admin confiável & Implementação RFC~3161 TSA; ``à prova de adulteração''; modelo de ameaças com 3 classes \\
\hline
R1.2 / R2.2 & Sem estatística inferencial & ICs de 95\%, Mann--Whitney $U$, Fisher exato, tamanhos de efeito, linguagem moderada \\
\hline
R1.3 / R2.1 & Novidade superestimada & Contribuição reposicionada como engenharia de integração \\
\hline
R1.4 / R2.5 & Comparação superficial & Tabela~3 redesenhada: escala de 3 níveis, 10 colunas, notas metodológicas \\
\hline
R1.5 & Categoria incorreta & Reclassificada como ``Ciência da Computação'' \\
\hline
R2.3 & Dados retrospectivos & Viés de memória e confundidores reconhecidos; subseções de limitações adicionadas \\
\hline
R2.6 & Sem validação de segurança & Verificação OWASP Top~10, npm audit; limitação honestamente declarada \\
\hline
R2.7 & Sem teste de estresse & Testes de carga executados: 1.753~req/s GET, 1.338~req/s POST @100 conn; saturação a 500 conn; JSON bruto no pacote de replicação \\
\hline
R2.8 & LGPD não demonstrada & Tabela de mapeamento LGPD, terminologia ``compatível com a LGPD'', disclaimer RIPD \\
\hline
R2.9 & Sem pacote de replicação & Pacote de replicação publicado no Zenodo (DOI: 10.5281/zenodo.18868994) \\
\hline
R2.10 & Prolixo & Texto condensado em $\sim$15\%; redundâncias removidas \\
\hline
\end{tabular}
\end{table}

\bigskip

\noindent Confiamos que estas revisões abordam adequadamente todas as preocupações levantadas por ambos os revisores e pelo editor. Permanecemos à disposição dos revisores para quaisquer esclarecimentos adicionais.

\bigskip
\noindent Atenciosamente,

\medskip
\noindent L.D.V.~Santos, C.V.S.~Oliveira, D.P.V.~Santos, T.C.~Azevedo,\\
R.N.~Araújo Filho, J.M.~Santos, M.M.~Fernandes

\end{document}
